\documentclass[11pt]{article}
\usepackage[T1]{fontenc}
\usepackage{amsmath,amssymb}
\usepackage{cite}
\usepackage{booktabs}
\usepackage{graphicx}
\usepackage{xurl}
\usepackage{hyperref}
\hypersetup{
  colorlinks=true,
  linkcolor=blue,
  urlcolor=blue,
  citecolor=blue,
  pdfauthor={Jabulani Monchwe, Martin Chanza, Olorato Motitswane, Mfundo Monchwe}
}
\newcommand{\software}{\textsf{ReliableML}}
\newcommand{\code}[1]{\texttt{\small #1}}
\newcommand{\jabulaniauthor}{
\IEEEauthorblockN{Jabulani Monchwe}
\IEEEauthorblockA{\textit{Applied Business Analytics and Decision Making}\\
\textit{North-West University}\\
\textit{IDEA-Research Club Chairperson}\\
Johannesburg, South Africa\\
\href{mailto:monchwej.25@gmail.com}{monchwej.25@gmail.com}}
}
\newcommand{\martinchanzauthor}{
\IEEEauthorblockN{Martin Chanza}
\IEEEauthorblockA{\textit{Applied Business Analytics and Decision Making}\\
\textit{North-West University}\\
Mafikeng, South Africa\\
\href{mailto:martin.chanza@nwu.ac.za}{martin.chanza@nwu.ac.za}}
}
\newcommand{\oloratoauthor}{
\IEEEauthorblockN{Olorato Motitswane}
\IEEEauthorblockA{\textit{Eduvos Midrand Campus -- IT Faculty Data Science, South Africa}\\
\textit{IDEA-Research Club}\\
\href{mailto:oloratomotitswane@gmail.com}{oloratomotitswane@gmail.com}}
}
\newcommand{\mfundoauthor}{
\IEEEauthorblockN{Mfundo Monchwe}
\IEEEauthorblockA{\textit{Eduvos Midrand Campus -- IT Faculty Data Science, South Africa}\\
\textit{IDEA-Research Club}\\
\href{mailto:diditmfundo@gmail.com}{diditmfundo@gmail.com}}
}
\newcommand{\minenhleauthor}{
\IEEEauthorblockN{Minenhle Mchunu}
\IEEEauthorblockA{\textit{Applied Business Analytics and Decision Making}\\
\textit{Eduvos Midrand Campus -- IT Faculty Data Science, South Africa}\\
\textit{IDEA-Research Club}\\
\href{mailto:Mchunu.minenhle02@gmail.com}{Mchunu.minenhle02@gmail.com}}
}
\newcommand{\projecturl}{\url{https://github.com/Jabualani/ML-Reliabilty}}

\title{\software{} Method Note: Group Fairness Metrics}
\author{\contactauthors}
\date{Version 0.1.1}
\begin{document}
\maketitle
\begin{abstract}
The fairness assessment compares observed classification rates across
user-supplied protected groups. It reports statistical parity,
equalized-odds, and predictive-parity gaps with denominators and bootstrap
intervals. The metrics are descriptive and do not select legal or
application-specific thresholds.
\end{abstract}

\section{Inputs}
Protected attributes are passed separately from model features to reduce the
risk of accidentally training on them. True labels and model predictions
must be row-aligned. Binary classification requires the user to specify the
positive label. Multiclass evaluation uses one-vs-rest calculations for each
observed/model class.

\section{Rate comparisons}
For each group, the assessment reports:
\begin{description}
  \item[Statistical parity:] the positive prediction or selection rate,
  \(P(\widehat{Y}=1\mid A=a)\), and maximum-minus-minimum gap across groups.
  \item[Equalized odds:] true-positive rates and false-positive rates
  conditional on the true class, with separate max-minus-minimum group gaps.
  \item[Predictive parity:] positive predictive value,
  \(P(Y=1\mid \widehat{Y}=1,A=a)\), and the corresponding group gap.
\end{description}
Group denominators are essential: an undefined rate caused by no eligible
observations is not interpreted as zero. Bootstrap intervals reflect
resampling variation under the implemented independent-row bootstrap.

\section{Interpretation}
A zero observed gap means rates were equal in the supplied sample under the
specified definitions; it is not proof of fairness, absence of harm, or
compliance with law. Estimates can be unstable for small subgroups or rare
outcomes. The three criteria can conflict, particularly when base rates
differ. Examine uncertainty, sample composition, intersectional groups,
measurement quality, and application-specific constraints before acting.

\section{Interface}
\begin{verbatim}
evaluation = model.evaluate(
    X_test, y_test,
    sensitive_features={"group": protected_test},
    positive_label=1,
)
print(evaluation.fairness)
\end{verbatim}

\section{Further reading}
\begin{thebibliography}{2}
\bibitem{hardt}
M. Hardt, E. Price, and N. Srebro. Equality of opportunity in supervised
learning. \emph{Advances in Neural Information Processing Systems}, 2016.
\url{https://arxiv.org/abs/1610.02413}.
\bibitem{reliability}
\software{} definitions and bootstrap reporting:
\url{https://github.com/Jabualani/ML-Reliabilty}.
\end{thebibliography}
\end{document}
